% MANUSCRIPT.tex, the manuscript skeleton for Paper A (Stereotype Rigidity in
% Social Inferences). Created 14 October 2026 as pass 1 of the plan in TODO.md
% (item PLAN), following papers/MANUSCRIPT_SKELETON.md. The predecessor
% document is papers/subjective_racism/theory_and_decisions/manuscript.tex,
% which stays in place until a later pass retires it.
%
% STRUCTURE THE CHECKER RELIES ON (checks/verify_manuscript.py)
%   Rows are written only through the four row macros below, and proofs only
%   through the mproof environment. IDs are K1, K2, ... (introduction),
%   M1, ... (model), L1, ... (literature) and C1, ... (conclusions).
%     \krow{ID}{headline}{model rows}{scope note}
%     \mrow{ID}{claim in words}{statement in symbols}{Lean theorems}
%     \lrow{ID}{bib key}{what we rely on}{verbatim quote}{page}{model rows}
%     \crow{ID}{statement}{model or literature rows}
%     \begin{mproof}{model row}{Lean theorems} ... \end{mproof}
%   Lists of rows or theorems are comma-separated. Lean theorem names are the
%   fully qualified names declared under lean/ or lean/mathlib/.
%
% RULES THE CHECKER ENFORCES
%   Introduction rows point to model rows only. Conclusion rows point to model
%   or literature rows only. Every model row names Lean theorems that exist in
%   an audited Lean file and has exactly one proof in Appendix A naming the
%   same theorems. Every literature row's quote appears verbatim in the
%   paper's lit/ record. In every proof, plain English is at most 30% of the
%   text, measured as prose words over prose words plus mathematical tokens.

\documentclass[11pt]{article}
\usepackage[T1]{fontenc}
\usepackage[utf8]{inputenc}
\usepackage{mathptmx}
\usepackage{microtype}
\usepackage[a4paper,margin=1.8cm]{geometry}
\usepackage{amsmath,amssymb,amsthm}
\usepackage{booktabs}
\usepackage{array}
\usepackage{longtable}
\usepackage{url}
\usepackage[hidelinks]{hyperref}
\usepackage{natbib}

\emergencystretch=1.5em

\newcolumntype{P}[1]{>{\raggedright\arraybackslash}p{#1}}

\urlstyle{tt}
\makeatletter
\g@addto@macro\UrlBreaks{\do\.\do\_\do\,}
\makeatother
\newcommand{\lean}[1]{\path{#1}}
\newcommand{\rowref}[1]{#1}
\newcommand{\unv}[1]{\textit{Unverified:} #1}
\newcommand{\pend}[1]{\textit{Pending:} #1}

\newcommand{\krow}[4]{#1 & #2 & \rowref{#3} & #4 \\ \addlinespace}
\newcommand{\mrow}[4]{#1 & #2 & #3 & \lean{#4} & p.~\pageref{proof:#1} \\ \addlinespace}
\newcommand{\lrow}[6]{#1 & \citet{#2} & #3 & ``#4'' & #5 & \rowref{#6} \\ \addlinespace}
\newcommand{\crow}[3]{#1 & #2 & \rowref{#3} \\ \addlinespace}

\newenvironment{mproof}[2]
  {\subsection*{Proof of #1}\label{proof:#1}%
   {\raggedright\noindent\textit{Lean:} \lean{#2}.\par}\smallskip}
  {\par\hfill$\square$\par\medskip}

\title{Stereotype Rigidity in Social Inferences}
\author{Anurag Srivastava}
\date{Manuscript skeleton, 14 October 2026}

\begin{document}
\maketitle

\section{Introduction}

% The question below recasts the first sentence of the predecessor's abstract.
% The author has not yet confirmed it (TODO.md, item Q).
Why do social boundaries persist despite counter-examples, legal equality and
anti-discrimination instruments?

\footnotesize
\begin{longtable}{@{}P{0.8cm}P{7.0cm}P{1.9cm}P{6.4cm}@{}}
\caption{Headline claims, each traced to the model rows it rests on.}\\
\toprule
\textbf{ID} & \textbf{Headline} & \textbf{Model rows} & \textbf{Scope note} \\
\midrule
\endfirsthead
\toprule
\textbf{ID} & \textbf{Headline} & \textbf{Model rows} & \textbf{Scope note} \\
\midrule
\endhead
\bottomrule
\endfoot
\end{longtable}
\normalsize

\section{Model}

\footnotesize
\begin{longtable}{@{}P{0.8cm}P{3.8cm}P{6.3cm}P{3.7cm}P{1.1cm}@{}}
\caption{Results of the model, each with its statement and the Lean theorems
that check it.}\\
\toprule
\textbf{ID} & \textbf{Claim} & \textbf{Statement} & \textbf{Lean} & \textbf{Proof} \\
\midrule
\endfirsthead
\toprule
\textbf{ID} & \textbf{Claim} & \textbf{Statement} & \textbf{Lean} & \textbf{Proof} \\
\midrule
\endhead
\bottomrule
\endfoot
\mrow{M1}{A marker whose likelihood does not depend on class leaves the
  Bayesian posterior on class equal to the prior. Under a prior that gives
  every class positive probability, the converse holds, so the two forms of
  group sufficiency are equivalent.}
  {$\mathcal{C}$ finite, $\sum_{c}\pi_g(c)=1$,
  $Z_g=\sum_{c'}\ell_g(c')\,\pi_g(c')$,
  $\pi'_g(c)=\ell_g(c)\,\pi_g(c)/Z_g$.
  (i) If $\ell_g(c)=k\neq0$ for every $c$, then $\pi'_g=\pi_g$.
  (ii) If $\pi_g(c)>0$ for every $c$, $Z_g\neq0$ and $\pi'_g=\pi_g$, then
  $\ell_g(c)=Z_g$ for every $c$.
  Neither hypothesis can be dropped.}
  {Sufficiency.posterior_eq_prior_of_flat, Sufficiency.flat_of_posterior_eq_prior, Sufficiency.control_flat_needed, Sufficiency.control_support_needed}
\mrow{M2}{Under group sufficiency, the Bayesian posterior on class is the
  within-group class distributions mixed by the Bayesian posterior on group.}
  {$\mathcal{G}$, $\mathcal{C}$ finite, $\sum_c\pi_g(c)=1$ and
  $\ell_g(c)=\lambda_g$ for every $g$ and $c$. Then
  $P(c\mid m)=\sum_g\pi_g(c)\,q'(g)$ for every $c$.
  It fails without $\ell_g(c)=\lambda_g$, and without
  $\sum_c\pi_g(c)=1$.}
  {ClassPosterior.bayes_class_eq_mixture, ClassPosterior.control_sufficiency_needed, ClassPosterior.control_normalised_needed}
\mrow{M3}{Every revision of the group weights from a prior that gives every
  group positive weight is a Bayesian update, for a marker likelihood that
  depends on group alone. The posterior on class is then the mixture by the
  revised weights.}
  {$q(g)>0$, $\sum_g w(g)=1$, $\sum_c\pi_g(c)=1$, $\kappa\neq0$ and
  $\ell_g(c)=\kappa\,w(g)/q(g)$. Then $q'=w$ and
  $P(c\mid m)=\sum_g\pi_g(c)\,w(g)$. If also $0<\kappa\le q(g)$ and
  $0\le w(g)\le1$, then $0\le\ell_g(c)\le1$.
  It fails without $q(g)>0$, since $q=(1,0)$ gives $q'(1)=0$ for every
  $\ell$.}
  {ClassPosterior.bayes_group_of_ratio, ClassPosterior.bayes_class_of_ratio, ClassPosterior.ratioLik_mem_unit, ClassPosterior.control_prior_support_needed}
\mrow{M4}{The composed map sends standing 0 to $\mu$ for every
  discrimination-rate function. It fixes standing 1 exactly when $F(1)=0$.
  With $F(\bar c)=(1-\bar c)^2$ it is the canonical map $\mathcal{H}$.}
  {$\mathcal{H}_F(0)=\mu$. If $s_g\delta\neq0$, then
  $\mathcal{H}_F(1)=1\iff F(1)=0$. For $F(\bar c)=(1-\bar c)^2$,
  $\mathcal{H}_F=\mathcal{H}$ with $b=s_g\delta$.
  It fails without $F(1)=0$, since $F\equiv1$, $\mu=\tfrac1{20}$,
  $\tilde\gamma=\tfrac1{10}$ and $s_g\delta=1$ give $\mathcal{H}_F(1)=0$.}
  {ClosureMap.Hf_zero, ClosureMap.Hf_one_iff, ClosureMap.H_eq_Hf, ClosureMap.control_F_one_needed}
\mrow{M5}{For the canonical map, there are two fixed points strictly inside
  $(0,1)$ exactly when the trap condition holds. The fixed points in $[0,1]$
  are then those two and integration at 1.}
  {$\mu>0$, $\tilde\gamma\ge0$, $b>0$.
  $\mathcal{H}(\bar c)-\bar c=(1-\bar c)\,Q(\bar c)$.
  $\exists\,0<c_L<c_M<1$ with $\mathcal{H}(c_L)=c_L$, $\mathcal{H}(c_M)=c_M$
  $\iff$ $b-\tilde\gamma>0$ and $(b-\tilde\gamma)^2>4\mu b$.
  Then $\mathcal{H}(\bar c)=\bar c\iff\bar c\in\{c_L,c_M,1\}$.
  The condition holds at $(\mu,\tilde\gamma,b)=(\tfrac1{20},\tfrac1{10},\tfrac12)$
  and fails at $(\tfrac1{20},\tfrac1{10},\tfrac15)$, where no fixed point lies
  inside $(0,1)$.}
  {ClosureMap.H_sub_self, ClosureMap.fixed_iff_Q, ClosureMap.H_one, ClosureMap.vieta, ClosureMap.Q_factor, ClosureMap.two_interior_iff, ClosureMap.fixed_points_exactly, ClosureMap.control_no_trap, ClosureMap.trap_at_figure_parameters}
\mrow{M6}{When the map is also increasing, every standing below the middle
  fixed point converges to the low one, the closure trap, and every standing
  above it converges to integration.}
  {$b>0$, $b<1-\mu+\tilde\gamma$, $\mu+\tilde\gamma<1$, and
  $0<c_L<c_M<1$ fixed points. Then $0\le x<y\le1\Rightarrow
  \mathcal{H}(x)<\mathcal{H}(y)$,
  $\bar c^0\in[0,c_M)\Rightarrow\mathcal{H}^n(\bar c^0)\to c_L$ and
  $\bar c^0\in(c_M,1]\Rightarrow\mathcal{H}^n(\bar c^0)\to1$.
  Both bounds hold at $(\tfrac1{20},\tfrac1{10},\tfrac12)$. Neither is
  automatic. At $(\tfrac1{20},0,3)$, $\mathcal{H}(\tfrac15)<0$, and at
  $(\tfrac1{20},0,1)$, $\mathcal{H}(\tfrac1{100})<\mathcal{H}(0)$.}
  {ClosureMap.tendsto_up, ClosureMap.tendsto_down, ClosureMap.H_continuous, ClosureMap.H_strictMonoOn, ClosureMap.basins, ClosureMap.basins_hypotheses_at_figure, ClosureMap.control_leaves_unit, ClosureMap.control_not_monotone}
\end{longtable}
\normalsize

\section{Literature}

\footnotesize
\begin{longtable}{@{}P{0.8cm}P{2.2cm}P{3.5cm}P{5.5cm}P{1.3cm}P{1.9cm}@{}}
\caption{Papers read, with the passage each row relies on.}\\
\toprule
\textbf{ID} & \textbf{Paper} & \textbf{What we rely on} & \textbf{Quote} &
\textbf{Page} & \textbf{Rows} \\
\midrule
\endfirsthead
\toprule
\textbf{ID} & \textbf{Paper} & \textbf{What we rely on} & \textbf{Quote} &
\textbf{Page} & \textbf{Rows} \\
\midrule
\endhead
\bottomrule
\endfoot
\end{longtable}
\normalsize

No paper has been read for this skeleton yet, so every claim the predecessor
manuscript makes about a source is unverified.

\section{Conclusions}

Table~\ref{tab:conclusions} will hold the conclusions that rest on model rows
and on papers read in full. Table~\ref{tab:pending} will hold every claim that
a result is new, and every other claim still waiting on a paper or a search.
Both are empty until the model rows they need exist.

\footnotesize
\begin{longtable}{@{}P{0.8cm}P{12.3cm}P{3.4cm}@{}}
\caption{Conclusions resting on model rows and papers read.}
\label{tab:conclusions}\\
\toprule
\textbf{ID} & \textbf{Conclusion} & \textbf{Rows} \\
\midrule
\endfirsthead
\toprule
\textbf{ID} & \textbf{Conclusion} & \textbf{Rows} \\
\midrule
\endhead
\bottomrule
\endfoot
\end{longtable}

\begin{longtable}{@{}P{0.8cm}P{12.3cm}P{3.4cm}@{}}
\caption{Pending conclusions, each naming what it waits on.}
\label{tab:pending}\\
\toprule
\textbf{ID} & \textbf{Conclusion} & \textbf{Rows} \\
\midrule
\endfirsthead
\toprule
\textbf{ID} & \textbf{Conclusion} & \textbf{Rows} \\
\midrule
\endhead
\bottomrule
\endfoot
\end{longtable}
\normalsize

\appendix

\section{Proofs}
\label{app:proofs}

\paragraph{Notation.} A group is $g$. Economic class takes values in a finite
set $\mathcal{C}$. The prior $\pi_g(c)=P(c\mid g)$ is the observer's
distribution of class within group $g$. The marker is $m$, and its likelihood
is $\ell_g(c)=P(m\mid c,g)$. The marginal of the marker within the group is
$Z_g=\sum_{c'\in\mathcal{C}}\ell_g(c')\,\pi_g(c')=P(m\mid g)$, and the
Bayesian posterior on class is $\pi'_g(c)=\ell_g(c)\,\pi_g(c)/Z_g=P(c\mid m,g)$.
Rows M2 and M3 range over a finite set of groups $\mathcal{G}$. The
observer's weights on the groups are $q(g)$ before the marker is read, the
marginal of the marker is
$Z=\sum_{g}\sum_{c}\ell_g(c)\,\pi_g(c)\,q(g)=P(m)$, the Bayesian posterior on
group is $q'(g)=\sum_{c}\ell_g(c)\,\pi_g(c)\,q(g)/Z$, and the Bayesian
posterior on class is $P(c\mid m)=\sum_{g}\ell_g(c)\,\pi_g(c)\,q(g)/Z$. In M3,
$w$ is a target revision of the group weights and $\kappa$ a scale.

Rows M4 to M6 concern a group's standing $\bar c\in[0,1]$, the mean class of
the group, which the observer's prior re-learns each generation
(Appendix~\ref{app:scope}). The parameters are the convergence rate $\mu$,
the mobility $\tilde\gamma$, the salience $s_g$, the cost of exclusion $\delta$
and the discrimination-rate function $F$. The composed map is
$\mathcal{H}_F(\bar c)=\mu+\bar c\,(1-\mu+\tilde\gamma-s_g\delta F(\bar c))
-\tilde\gamma\,\bar c^2$. With $F(\bar c)=(1-\bar c)^2$ and $b=s_g\delta$ it is
$\mathcal{H}(\bar c)=\mu+\bar c\,(1-\mu+\tilde\gamma-b(1-\bar c)^2)
-\tilde\gamma\,\bar c^2$, and $Q(\bar c)=b\bar c^2-(b-\tilde\gamma)\bar c+\mu$.
$\mathcal{H}^n$ is the $n$-th iterate of $\mathcal{H}$.

\begin{mproof}{M1}{Sufficiency.posterior_eq_prior_of_flat, Sufficiency.flat_of_posterior_eq_prior, Sufficiency.control_flat_needed, Sufficiency.control_support_needed}
(i) With $\ell_g(c)=k\neq0$ for every $c\in\mathcal{C}$ and $\sum_c\pi_g(c)=1$,
\[
Z_g=\sum_{c'}k\,\pi_g(c')=k\sum_{c'}\pi_g(c')=k,
\qquad
\pi'_g(c)=\frac{k\,\pi_g(c)}{k}=\pi_g(c).
\]
(ii) With $Z_g\neq0$, $\pi_g(c)>0$ and $\pi'_g(c)=\pi_g(c)$,
\[
\ell_g(c)\,\pi_g(c)=Z_g\,\pi_g(c)
\;\Longrightarrow\;
\ell_g(c)=Z_g .
\]
Controls, with $\mathcal{C}=\{0,1\}$. For (i), $\pi_g=(\tfrac12,\tfrac12)$ and
$\ell_g=(1,0)$ give $Z_g=\tfrac12$ and $\pi'_g(0)=1\neq\tfrac12=\pi_g(0)$. For
(ii), $\pi_g=(1,0)$ and $\ell_g=(1,2)$ give $Z_g=1$ and $\pi'_g=\pi_g$, while
$\ell_g(0)=1\neq2=\ell_g(1)$.
\end{mproof}

\begin{mproof}{M2}{ClassPosterior.bayes_class_eq_mixture, ClassPosterior.control_sufficiency_needed, ClassPosterior.control_normalised_needed}
With $\ell_g(c)=\lambda_g$ and $\sum_c\pi_g(c)=1$,
\[
\sum_{c'}\ell_g(c')\,\pi_g(c')\,q(g)=\lambda_g\,q(g)\sum_{c'}\pi_g(c')=\lambda_g\,q(g)
\;\Longrightarrow\;
q'(g)=\frac{\lambda_g\,q(g)}{Z},
\]
\[
P(c\mid m)=\sum_g\frac{\lambda_g\,\pi_g(c)\,q(g)}{Z}
=\sum_g\pi_g(c)\,\frac{\lambda_g\,q(g)}{Z}=\sum_g\pi_g(c)\,q'(g).
\]
Controls, with $\mathcal{G}=\{0\}$ and $\mathcal{C}=\{0,1\}$. For sufficiency,
$q=(1)$, $\pi_0=(\tfrac12,\tfrac12)$ and $\ell_0=(1,0)$ give $Z=\tfrac12$,
$P(0\mid m)=1$ and $q'(0)=1$, so $\sum_g\pi_g(0)\,q'(g)=\tfrac12\neq1$. For
normalisation, $\pi_0=(1,1)$ and $\ell_0=(1,1)$ give $Z=2$, $P(0\mid m)=\tfrac12$
and $q'(0)=1$, so $\sum_g\pi_g(0)\,q'(g)=1\neq\tfrac12$.
\end{mproof}

\begin{mproof}{M3}{ClassPosterior.bayes_group_of_ratio, ClassPosterior.bayes_class_of_ratio, ClassPosterior.ratioLik_mem_unit, ClassPosterior.control_prior_support_needed}
With $\ell_g(c)=\kappa\,w(g)/q(g)$, $q(g)>0$ and $\sum_c\pi_g(c)=1$,
\[
\sum_c\ell_g(c)\,\pi_g(c)\,q(g)=\kappa\,w(g)\sum_c\pi_g(c)=\kappa\,w(g),
\qquad
Z=\kappa\sum_g w(g)=\kappa,
\qquad
q'(g)=\frac{\kappa\,w(g)}{\kappa}=w(g).
\]
$\ell_g$ is constant in $c$, so M2 applies and
$P(c\mid m)=\sum_g\pi_g(c)\,q'(g)=\sum_g\pi_g(c)\,w(g)$. With
$0<\kappa\le q(g)$ and $0\le w(g)\le1$,
\[
0\le\frac{\kappa\,w(g)}{q(g)}
\quad\text{and}\quad
\kappa\,w(g)\le\kappa\le q(g)\;\Longrightarrow\;\frac{\kappa\,w(g)}{q(g)}\le1.
\]
Control, with $\mathcal{G}=\{0,1\}$ and $q=(1,0)$. For every $\ell$ and $\pi$,
$\sum_c\ell_1(c)\,\pi_1(c)\cdot0=0$, so $q'(1)=0$.
\end{mproof}

\begin{mproof}{M4}{ClosureMap.Hf_zero, ClosureMap.Hf_one_iff, ClosureMap.H_eq_Hf, ClosureMap.control_F_one_needed}
\[
\mathcal{H}_F(0)=\mu+0-0=\mu,
\qquad
\mathcal{H}_F(1)=\mu+(1-\mu+\tilde\gamma-s_g\delta F(1))-\tilde\gamma
=1-s_g\delta F(1).
\]
So $\mathcal{H}_F(1)=1\iff s_g\delta F(1)=0\iff F(1)=0$ when $s_g\delta\neq0$.
Substituting $F(\bar c)=(1-\bar c)^2$ and $b=s_g\delta$ gives $\mathcal{H}$.
Control: $F\equiv1$ gives
$\mathcal{H}_F(1)=\tfrac1{20}+(1-\tfrac1{20}+\tfrac1{10}-1)-\tfrac1{10}=0$.
\end{mproof}

\begin{mproof}{M5}{ClosureMap.H_sub_self, ClosureMap.fixed_iff_Q, ClosureMap.H_one, ClosureMap.vieta, ClosureMap.Q_factor, ClosureMap.two_interior_iff, ClosureMap.fixed_points_exactly, ClosureMap.control_no_trap, ClosureMap.trap_at_figure_parameters}
Expanding,
\[
\mathcal{H}(\bar c)-\bar c=\mu(1-\bar c)+\tilde\gamma\,\bar c(1-\bar c)
-b\,\bar c(1-\bar c)^2=(1-\bar c)\,Q(\bar c),
\]
so for $\bar c\neq1$, $\mathcal{H}(\bar c)=\bar c\iff Q(\bar c)=0$, and
$\mathcal{H}(1)=1$.
($\Rightarrow$) $Q(c_L)=Q(c_M)=0$ with $c_L\neq c_M$ gives
\[
b(c_L+c_M)=b-\tilde\gamma,\qquad b\,c_Lc_M=\mu,
\]
so $b-\tilde\gamma=b(c_L+c_M)>0$ and
$(b-\tilde\gamma)^2-4\mu b=b^2(c_L+c_M)^2-4b^2c_Lc_M=b^2(c_M-c_L)^2>0$.
($\Leftarrow$) Put $B=b-\tilde\gamma$, $D=B^2-4\mu b>0$, $r=\sqrt D$ and
\[
c_L=\frac{B-r}{2b},\qquad c_M=\frac{B+r}{2b}.
\]
Then $Q(c_L)=Q(c_M)=0$ by $r^2=D$. $0<r<B$ since $D<B^2$, so $0<c_L<c_M$.
$r<b+\tilde\gamma$ since $D-(b+\tilde\gamma)^2=-4b\tilde\gamma-4\mu b<0$, so
$c_M<1$.
Exactness. With the two relations above,
$Q(\bar c)=b(\bar c-c_L)(\bar c-c_M)$, so
\[
\mathcal{H}(\bar c)-\bar c=(1-\bar c)\,b\,(\bar c-c_L)(\bar c-c_M)=0
\iff\bar c\in\{c_L,c_M,1\}.
\]
Controls. At $(\tfrac1{20},\tfrac1{10},\tfrac12)$, $B^2=\tfrac4{25}>\tfrac1{10}=4\mu b$.
At $(\tfrac1{20},\tfrac1{10},\tfrac15)$, $B^2=\tfrac1{100}<\tfrac1{25}=4\mu b$.
\end{mproof}

\begin{mproof}{M6}{ClosureMap.tendsto_up, ClosureMap.tendsto_down, ClosureMap.H_continuous, ClosureMap.H_strictMonoOn, ClosureMap.basins, ClosureMap.basins_hypotheses_at_figure, ClosureMap.control_leaves_unit, ClosureMap.control_not_monotone}
Put $A=1-\mu+\tilde\gamma-b$. For $0\le x<y\le1$,
\[
\mathcal{H}(y)-\mathcal{H}(x)=(y-x)\,G,\qquad
G=A+(2b-\tilde\gamma)(x+y)-b(x^2+xy+y^2),
\]
\[
G=A(1-x)(1-y)+(1-\mu)\big(x(1-y)+y(1-x)\big)+(1-\mu-\tilde\gamma)\,xy
+b\,x(1-x)+b\,y(1-y).
\]
The three weights are $\ge0$ and sum to 1, $A>0$, $1-\mu>0$ and
$1-\mu-\tilde\gamma>0$, so $G\ge\min(A,1-\mu-\tilde\gamma)>0$.
By M5, $\mathcal{H}(\bar c)-\bar c=(1-\bar c)\,b\,(\bar c-c_L)(\bar c-c_M)$, which is
$>0$ on $[0,c_L)$, $<0$ on $(c_L,c_M)$ and $>0$ on $(c_M,1)$.
Monotone orbits. Let $f$ be continuous, nondecreasing on $[x_0,p]$, $f(p)=p$
and $f(x)>x$ on $[x_0,p)$. With $u_n=f^n(x_0)$,
\[
x_0\le u_n\le u_{n+1}\le p,\qquad u_n\to L=\sup_n u_n\in[x_0,p],
\qquad f(L)=\lim f(u_n)=\lim u_{n+1}=L,
\]
and $L<p$ would give $f(L)>L$, so $L=p$. The mirror case $f(x)<x$ on
$(p,x_0]$ gives $u_n\downarrow p$.
Apply it with $p=c_L$ for $\bar c^0\in[0,c_L]$, the mirror with $p=c_L$ for
$\bar c^0\in(c_L,c_M)$, and $p=1$ for $\bar c^0\in(c_M,1]$, where
$\mathcal{H}(1)=1$ by M5.
Controls. $\mathcal{H}(\tfrac15)=\tfrac1{20}+\tfrac15(\tfrac{19}{20}-3\cdot\tfrac{16}{25})<0$
at $(\tfrac1{20},0,3)$, and
$\mathcal{H}(\tfrac1{100})-\mathcal{H}(0)=\tfrac1{100}(\tfrac{19}{20}-(\tfrac{99}{100})^2)<0$
at $(\tfrac1{20},0,1)$.
\end{mproof}

\section{Evidence and verification}
\label{app:evidence}

Every row of the model table is a theorem of Lean~4 with Mathlib. The Lean
column names the theorems, and the proof of each row in
Appendix~\ref{app:proofs} names the same ones. The development lives in
\path{lean/mathlib/}, built on the toolchain \path{v4.24.0-rc1} against
Mathlib commit \path{14871d5}. The script \path{verify.sh} runs three suites on
every build, and \path{VERIFICATION.md} records what they printed.

\begin{enumerate}
\item \emph{The Lean build and axiom audit.} Every file under
\path{lean/mathlib/} is a build root and builds. No source contains
\texttt{sorry}. Every declared theorem is audited with \texttt{\#print axioms},
and none depends on an axiom outside \texttt{propext},
\texttt{Classical.choice} and \texttt{Quot.sound}.
\item \emph{The skeleton checker.} Rules MS-1 to MS-6 of
\path{checks/verify_manuscript.py}, each also run on a well-formed sample it
must accept and a faulty sample it must reject.
\item \emph{Doc consistency.} Every theorem count stated here is in digits and
matches the audit.
\end{enumerate}

\paragraph{What Lean does not carry.} The class space is finite in the Lean
statements of M1 to M3, and so is the set of groups in M2 and M3. Rows M5
and M6 are proved for the canonical $F(\bar c)=(1-\bar c)^2$ only. The general
$F$ of the predecessor's regularity assumption is not yet covered. The
composed map itself enters every row as a definition, so its derivation from
the engagement decision is not checked (Appendix~\ref{app:scope}). The predecessor manuscript gives class a Beta law on
$[0,1]$, and the continuous form of M1 is not formalised.

\paragraph{Illustrations.} No row rests on SymPy or on numerical sampling. The
SymPy scripts of the predecessor bundle are not part of this skeleton.

\section{Terminology}
\label{app:terms}

\footnotesize
\begin{longtable}{@{}P{3.4cm}P{4.2cm}P{8.6cm}@{}}
\toprule
\textbf{Term} & \textbf{Symbol} & \textbf{Meaning} \\
\midrule
\endfirsthead
\toprule
\textbf{Term} & \textbf{Symbol} & \textbf{Meaning} \\
\midrule
\endhead
\bottomrule
\endfoot
Prior & $\pi_g(c)$ & The observer's distribution of class within group $g$
  before the marker is read. \\ \addlinespace
Likelihood & $\ell_g(c)$ & The probability of the marker given class and
  group. \\ \addlinespace
Marginal & $Z_g$ & The probability of the marker given the group alone. \\
  \addlinespace
Posterior & $\pi'_g(c)$ & The Bayesian distribution of class within group $g$
  after the marker is read. \\ \addlinespace
Group sufficiency & $\ell_g(c)=Z_g$ for every $c$ & The marker says nothing
  about class beyond what the group says. M1 shows this equals
  $\pi'_g=\pi_g$ under a prior with full support. \\ \addlinespace
Group weights & $q(g)$ & The observer's probability that the person
  read belongs to group $g$, before the marker is read. \\ \addlinespace
Posterior on group & $q'(g)$ & The Bayesian group weights after the marker
  is read. \\ \addlinespace
Mixture & $\sum_g\pi_g(c)\,w(g)$ & The distribution of class reached by
  weighting each group's class distribution by group weights $w$. M2 and M3
  show it is the Bayesian posterior on class exactly when $w=q'$. \\
  \addlinespace
Standing & $\bar c$ & The mean class of a group, in $[0,1]$. \\ \addlinespace
Composed map & $\mathcal{H}_F$, $\mathcal{H}$ & A group's standing in the
  next generation as a function of its standing now. \\ \addlinespace
Trap condition & $b-\tilde\gamma>0$, $(b-\tilde\gamma)^2>4\mu b$ & The
  condition of M5 for two fixed points inside $(0,1)$. \\ \addlinespace
Closure trap & $c_L$ & The low interior fixed point. M6 shows it attracts
  every standing below $c_M$. \\ \addlinespace
Tipping point & $c_M$ & The middle fixed point, which separates the two
  basins of M6. \\ \addlinespace
Integration & $\bar c=1$ & The fixed point at full standing (M4). \\
  \addlinespace
\end{longtable}
\normalsize

\section{Primitives and scope}
\label{app:scope}

\paragraph{Class space.} Class takes finitely many values in the Lean
statement of every row so far. Row M1 depends on this. The continuous class of
the predecessor manuscript lies outside the rows until a row states it.

\paragraph{Group sufficiency.} The predecessor manuscript states the
assumption in two forms joined by ``equivalently'', the likelihood form
$P(m\mid c,g)=P(m\mid g)$ and the posterior form $P(c\mid m,g)=P(c\mid g)$.
Row M1 proves the two forms equivalent when the prior gives every class
positive probability, and its second control shows the posterior form is
weaker without that support.

\paragraph{The updating rule.} The observer updates by Bayes' rule alone, a
choice the author made on 14 October 2026. The reason is rows M2 and M3. The
predecessor manuscript adds Jeffrey's rule of conditioning in the form of its
equation (kinematics), which mixes the within-group class distributions by
group weights taken as given. Under group sufficiency, the Bayesian posterior
on class is such a mixture (M2), and every such mixture whose prior gives each
group positive weight is a Bayesian posterior for some marker likelihood
(M3). Within one reading of a marker, the two rules therefore produce the same
set of posteriors on class, and a second rule adds no result. What lies
outside the rows is a sequence of several soft readings, where the order of
the readings could matter. The model has no such sequence.
\unv{The attribution of the kinematics formula to Jeffrey (1965) rests on the
predecessor manuscript. The source has not been read for this skeleton.}

\paragraph{The law of motion.} Rows M4 to M6 take the predecessor's
composed map, its equation (composed), as the law of motion of a group's
standing. Its derivation from the engagement decision and the discrimination
rate is not a row yet. The bounds $b<1-\mu+\tilde\gamma$ and
$\mu+\tilde\gamma<1$ of M6 limit the cost of exclusion from above. The
predecessor's ``sufficiently large $\delta$'' is therefore bounded, and the
first control of M6 shows the map leaves $[0,1]$ when $s_g\delta$ is large.

\paragraph{The group prior across generations.} Each generation the observer's
class distribution for a group is re-learned from the group's current
standing, a choice the author made on 14 October 2026. The reason is that
this is the model the predecessor manuscript writes down, in its equations
(prior), (dynamic) and (composed) and its loop figure, where the discrimination
rate of a generation is computed from that generation's standing. The
alternative holds the distribution fixed at its first value. No row depends on
this choice yet. The rows on the composed map and the closure trap will
(\path{TODO.md}, items Q3 to Q5), and the fixed alternative will enter as a
row of its own (Q5).

\section{Refuted conjectures}
\label{app:refuted}

None yet.

\section{Responses to the referee reports}

Not applicable. The paper has not been through review in this form.

\section{Coverage of the referee comments}

Not applicable, for the same reason.

\end{document}
